\section{Shape estimates and supporting functionals}
\label{df:sec:setup}

Fix $\kappa,\lambda>0$ throughout this part and write
$\mu=\mu_{\kappa,\lambda}$. All constants and scale thresholds may
depend on this fixed law. We shall compare path costs with the two
extreme supporting functionals at a prescribed direction. The crossing
estimates use only their one-sided contact with $\mu$; the assumption
that they differ enters the final contradiction.

\subsection{Uniform shape and local sums}

We first record the precise form of the standard shape theorem that
will be used. It holds for the fixed Gamma law:
\begin{equation}\label{df:eq:uniform-shape}
 \sup_{z\in\Z^2:\,|z|\ge R}
 \frac{|T(0,z)-\mu(z)|}{|z|}\longrightarrow0
 \qquad\text{almost surely as }R\to\infty.
\end{equation}
This is the shape theorem of
\cite[Theorem~A.1 and Equation~(A.3)]{KRS23}, at weight shift zero.
The Gamma law has essential infimum zero and no atom there. For
four independent copies of an edge weight,
\[
 \E\!\left[\min\{\tau_1,\tau_2,\tau_3,\tau_4\}^2\right]
 \le \E[\tau_1^2]
 =\frac{\kappa(\kappa+1)}{\lambda^2}<\infty.
\]
These facts verify the cited hypotheses. The source uses $\ell^1$
distance, which gives \eqref{df:eq:uniform-shape} by equivalence of
norms. There are deterministic constants $0<\alpha\le\beta<\infty$ such that
\begin{equation}\label{df:eq:norm-bounds}
 \alpha|z|\le\mu(z)\le\beta|z|\qquad(z\in\R^2).
\end{equation}
An event expressed relative to a deterministic lattice origin has the
same probability at every translate. This observation will allow us to
use shape estimates in probability at origins that vary with the
scale. Equation~\eqref{df:eq:uniform-shape} itself is an almost-sure
statement simultaneously over all endpoints from the fixed origin.

Almost surely every edge weight is finite and positive, since there
are only countably many edges. In particular, loop erasure does not
increase cost. A path constrained to a set will mean a path all of
whose vertices belong to that set; such path-existence events use
only its internal edges. Concatenated paths may retrace edges, with
each traversal contributing to the cost.

The following bound controls the vertical attachments used below.

\begin{lemma}[Gamma exponential moments]\label{df:ga:moment}
For $\tau\sim\operatorname{Gamma}(\kappa,\lambda)$ and $\theta<\lambda$,
\[
 \E e^{\theta\tau}
 =\left(\frac{\lambda}{\lambda-\theta}\right)^\kappa.
\]
If $V$ is the sum over a deterministic set of $M$ distinct edges and
$M\le A v$, where $A,v>0$, then
\begin{equation}\label{df:eq:gamma-sum-tail}
 \PP\left(V>\frac{2A\kappa}{\lambda}v\right)
 \le \exp\{-A\kappa(1-\log 2)v\}.
\end{equation}
\end{lemma}

\begin{proof}
Integrating the Gamma density with the substitution
$y=(\lambda-\theta)t$ gives the moment identity. At
$\theta=\lambda/2$ it reads $\E e^{\lambda\tau/2}=2^\kappa$.
Independence and exponential Markov inequality give
\[
 \PP\left(V>\frac{2A\kappa}{\lambda}v\right)
 \le e^{-A\kappa v}2^{\kappa M}
 \le e^{-A\kappa(1-\log 2)v}.
\]
\end{proof}

\subsection{The two extreme supports}

Signed coordinate permutations preserve the law and hence $\mu$.
Together with positive homogeneity, these symmetries reduce the proof
to directions
\[
 u=(1,b),\qquad -1\le b\le1.
\]
Fix such a deterministic direction for the rest of the argument, and
write
\[
 x(z)=z_1,\qquad r(z)=z_2-bz_1,
 \qquad m_0=\mu(u)>0.
\]
The function $t\mapsto\mu(u+te_2)$ is convex and Lipschitz. Let
\[
 a_- = \lim_{t\uparrow0}\frac{\mu(u+te_2)-m_0}{t},
 \qquad
 a_+ = \lim_{t\downarrow0}\frac{\mu(u+te_2)-m_0}{t}.
\]
These finite limits satisfy $a_-\le a_+$. Define
\begin{equation}\label{df:eq:support-definitions}
 \begin{split}
 f_\pm(z)&=m_0x(z)+a_\pm r(z),\\
 a&=\frac{a_-+a_+}{2},\qquad
 c=\frac{a_+-a_-}{2},\qquad
 f(z)=m_0x(z)+ar(z).
 \end{split}
\end{equation}

\begin{lemma}[One-sided contact]\label{df:prelim:supports}
The linear functionals $f_-$, $f_+$, and $f$ are bounded above by
$\mu$ on $\R^2$ and equal $\mu$ at $u$. For
$\sigma\in\{-,+\}$, with the corresponding numerical sign when
multiplied by a scalar,
\begin{equation}\label{df:eq:support-contact}
 \mu(u+te_2)-f_\sigma(u+te_2)=o(|t|)
 \qquad(t\to0,\ \sigma t>0).
\end{equation}
\end{lemma}

\begin{proof}
Each of $a_-$ and $a_+$ is a subgradient at zero of the convex
function $t\mapsto\mu(u+te_2)$. If $x(z)>0$, write
$z=x(z)(u+r(z)x(z)^{-1}e_2)$ and use positive homogeneity to get
$f_\sigma(z)\le\mu(z)$. If this inequality failed at any other
point $z$, then a sufficiently small positive convex combination of
$z$ with $u$ would still violate it, by convexity, and would have
positive first coordinate. This is impossible. The average $f$ is
also bounded above by $\mu$. Equation~\eqref{df:eq:support-contact}
is the respective one-sided derivative expansion.
\end{proof}

The \emph{defect} of a path is its cost minus the value of a supporting
functional on its displacement. We shall use
\begin{equation}\label{df:eq:defects}
 H(P)=\tau(P)-f(\Delta P),\qquad
 H_\sigma(P)=\tau(P)-f_\sigma(\Delta P).
\end{equation}
They are additive under concatenation, and
\begin{equation}\label{df:eq:minimum-defect}
 \min\{H_+(P),H_-(P)\}=H(P)-c|r(\Delta P)|.
\end{equation}
The crossing and endpoint estimates will hold for both signs whether
or not $c$ is positive. In Section~\ref{df:sec:witnesses}, the gap
assumption $c>0$ and \eqref{df:eq:minimum-defect} will make transverse
exits expensive relative to the average support $f$.
